\chapter{Bolt Pretension}
\label{chap:pretension}

Pretension in linear statics requires two coincident, planar solid faces with
separate node IDs. The preprocessor must split the mesh. FEMaster does not
modify topology. Reference-coordinate matching must be bijective: missing or
ambiguous matches are rejected.

\keywordpage{PRETENSIONSECTION}

\begin{keywordparameters}
\key{NAME} & required & Unique section name. \\
\key{SURFACE\_A} & required & First solid face set. \\
\key{SURFACE\_B} & required & Coincident, topologically separate second set. \\
\end{keywordparameters}

The section axis is the area-weighted normal of \val{SURFACE\_A}. For every
paired node, both tangential relative displacements vanish. All pairs share
one axial relative displacement in \val{FORCE} mode; this uses only physical
nodal degrees of freedom.

\keywordpage{PRETENSION}

\begin{keywordparameters}
\key{SECTION} & required & Previously defined section. \\
\key{MODE} & required & \val{FORCE} or \val{HOLD}. \\
\key{VALUE} & for FORCE & Total signed force on side B. \\
\end{keywordparameters}

The force is distributed equally across node pairs and applied in opposite
directions on the two faces. HOLD fixes the common axial gap from the preceding
successfully completed step, while subsequent loads may change the force.
Before its first action, a section has zero gap. No control node or additional
degree of freedom is created.

\begin{femasterdeck}
*PRETENSIONSECTION, NAME=PT1, SURFACE_A=CUT_A, SURFACE_B=CUT_B
*LOADCASE, TYPE=LINEARSTATIC, NAME=PRELOAD
*PRETENSION, SECTION=PT1, MODE=FORCE, VALUE=1000
*CONSTRAINTMETHOD, TYPE=NULLSPACE
*END
*LOADCASE, TYPE=LINEARSTATIC, NAME=SERVICE
*PRETENSION, SECTION=PT1, MODE=HOLD
*CONSTRAINTMETHOD, TYPE=NULLSPACE
*END
\end{femasterdeck}

LAGRANGE, nonlinear and perturbation steps are unsupported for pretension.

\section{Results}

\val{PTGAP} gives the axial relative displacement on both nodes of each pair.
\val{PTFORC} gives opposite transmitted force vectors on the two faces.
The signed section force is the sum of \val{PTFORC} on side B projected onto
the section axis. \val{TF} repeats that total on side B nodes. Supports at
the cut prevent unambiguous force recovery.